\documentclass[aspectratio=169,10pt]{beamer}

% ── Tema ──────────────────────────────────────────────────────────────────────
\usetheme{Madrid}
\usecolortheme{default}
\usefonttheme{professionalfonts}

% Paleta MendesLobo: azul estrutura + laranja acento
\definecolor{mlblue}{RGB}{18,52,86}
\definecolor{mlorange}{RGB}{214,100,28}
\definecolor{mlgray}{RGB}{240,242,245}
\definecolor{mlgreen}{RGB}{39,120,71}
\definecolor{mlred}{RGB}{180,35,35}

\setbeamercolor{structure}{fg=mlblue}
\setbeamercolor{frametitle}{bg=mlblue,fg=white}
\setbeamercolor{title}{bg=mlblue,fg=white}
\setbeamercolor{block title}{bg=mlblue,fg=white}
\setbeamercolor{block body}{bg=mlgray}
\setbeamercolor{alerted text}{fg=mlorange}
\setbeamercolor{example text}{fg=mlgreen}
\setbeamertemplate{navigation symbols}{}

% ── Pacotes ───────────────────────────────────────────────────────────────────
\usepackage[utf8]{inputenc}
\usepackage[T1]{fontenc}
\usepackage[portuguese]{babel}
\usepackage{amsmath,amssymb}
\usepackage{booktabs}
\usepackage{tabularx}
\usepackage{array}
\usepackage{multirow}
\usepackage{graphicx}
\usepackage{tikz}
\usetikzlibrary{arrows.meta,positioning,shapes.geometric,fit,backgrounds,calc}
\usepackage{listings}
\usepackage{xcolor}
\usepackage{fontawesome5}

\lstset{
  basicstyle=\ttfamily\scriptsize,
  breaklines=true,
  backgroundcolor=\color{mlgray},
  frame=single,
  rulecolor=\color{mlblue!30},
  xleftmargin=2pt,xrightmargin=2pt,
}

% ── Comandos utilitários ───────────────────────────────────────────────────────
\newcommand{\yes}{\textcolor{mlgreen}{\faCheck}}
\newcommand{\no}{\textcolor{mlred}{---}}
\newcommand{\half}{\textcolor{mlorange}{$\circ$}}
\newcommand{\role}[1]{\texttt{\small #1}}
\newcommand{\tool}[1]{\texttt{\footnotesize #1}}
\newcommand{\skill}[2]{\textbf{S#1}~#2}

% ── Metadados ─────────────────────────────────────────────────────────────────
\title[Concilium — Visão Técnica]{Concilium\\[4pt]\large Pipeline, Catálogo, Roles \& Workflows}
\subtitle{Mendes Lobo Construção}
\author{Bruno Peixoto}
\date{\today}
\institute{Uso interno --- confidencial}

% ── Início ────────────────────────────────────────────────────────────────────
\begin{document}

% ─────────────────────────────────────────────────────────────────────────────
\begin{frame}[plain]
  \titlepage
\end{frame}

% ─────────────────────────────────────────────────────────────────────────────
\begin{frame}{Agenda}
  \tableofcontents[hideallsubsections]
\end{frame}

% =============================================================================
\section{Arquitetura Geral}
% =============================================================================

\begin{frame}{Concilium em uma Página}
  \begin{columns}[T]
    \column{0.52\textwidth}
    \begin{block}{O que é}
      Sistema de \textbf{orçamentação e planejamento} de obras de infraestrutura
      que combina:
      \begin{itemize}
        \item Pipeline de skills determinística e auditável
        \item Catálogo versionado de insumos e composições (CPU)
        \item Fachada MCP para agentes de IA e usuários humanos
        \item Pool de recursos compartilhados com otimização RCPSP
      \end{itemize}
    \end{block}
    \vspace{4pt}
    \begin{block}{Fluxo de valor}
      {\small Lista de serviços $\;\to\;$ EAP $\;\to\;$ Cronograma $\;\to\;$
      Pool $\;\to\;$ \textbf{Resultado financeiro auditável}}
    \end{block}

    \column{0.44\textwidth}
    \begin{tikzpicture}[
        box/.style={draw=mlblue,fill=mlgray,rounded corners=3pt,
                    minimum width=3.2cm,minimum height=0.55cm,
                    font=\scriptsize\ttfamily,align=center},
        arr/.style={-{Stealth[length=5pt]},mlblue,thick},
        every node/.style={font=\scriptsize}
      ]
      \node[box] (mcp)   {Fachada MCP (48 tools)};
      \node[box,below=0.25cm of mcp] (pipe) {Pipeline de Skills (08)};
      \node[box,below=0.25cm of pipe] (can)  {Canônico JSON (versionado)};
      \node[box,below=0.25cm of can]  (cat)  {Catálogo (contexto)};
      \node[box,below=0.25cm of cat]  (infra){Azure ACA (blue-green)};
      \draw[arr] (mcp) -- (pipe);
      \draw[arr] (pipe) -- (can);
      \draw[arr] (can) -- (cat);
      \draw[arr] (cat) -- (infra);
    \end{tikzpicture}
  \end{columns}
\end{frame}

% =============================================================================
\section{Pipeline de Skills}
% =============================================================================

\begin{frame}{Pipeline: Composição de 8 Skills}
  \begin{tikzpicture}[
      skill/.style={draw=mlblue,fill=mlgray,rounded corners=3pt,
                    minimum width=2.1cm,minimum height=1.0cm,align=center,
                    font=\tiny},
      arr/.style={-{Stealth[length=4pt]},mlblue,thick},
      lbl/.style={font=\tiny\itshape,text=mlblue,align=center},
    ]
    % Nodes
    \node[skill] (s1) {\textbf{01}\\\scriptsize Escopo\\EAP};
    \node[skill,right=0.5cm of s1] (s2) {\textbf{02}\\\scriptsize Prod.\\Climática};
    \node[skill,right=0.5cm of s2] (s3) {\textbf{03}\\\scriptsize Catálogo\\CPU};
    \node[skill,right=0.5cm of s3] (s4) {\textbf{04}\\\scriptsize Cronograma\\RCPSP};
    \node[skill,right=0.5cm of s4] (s5) {\textbf{05}\\\scriptsize Equipes\\Pool};
    \node[skill,below=0.6cm of s5] (s6) {\textbf{06}\\\scriptsize Resultado\\Financeiro};
    \node[skill,left=0.5cm of s6]  (s7) {\textbf{07}\\\scriptsize Indiretos\\BDI};
    \node[skill,left=0.5cm of s7]  (s8) {\textbf{08}\\\scriptsize Artefatos\\Entrega};
    % Arrows horizontal
    \foreach \a/\b in {s1/s2,s2/s3,s3/s4,s4/s5}
      \draw[arr] (\a) -- (\b);
    % Turn-down
    \draw[arr] (s5) -- (s6);
    % Arrows return
    \foreach \a/\b in {s6/s7,s7/s8}
      \draw[arr] (\a) -- (\b);
    % Canonical label
    \node[lbl,above=0.15cm of s3] {Canônico JSON};
    \node[lbl,below=0.15cm of s7] {Canônico JSON};
  \end{tikzpicture}

  \vspace{4pt}
  \begin{columns}[T]
    \column{0.5\textwidth}
    \begin{itemize}\scriptsize
      \item Cada skill recebe o canônico da skill anterior como \textbf{entrada}
      \item Saída é persistida e versionada (imutável após aprovação)
      \item Execução assíncrona via worker pool (\texttt{disparar\_pipeline})
    \end{itemize}
    \column{0.5\textwidth}
    \begin{itemize}\scriptsize
      \item Skills são \textbf{idempotentes}: mesmo input $\to$ mesmo output
      \item Re-execução com novo catálogo ou override gera novo plano
      \item Artefatos (PDF/XLSX) gerados na skill 08
    \end{itemize}
  \end{columns}
\end{frame}

% ── Skills em detalhe ─────────────────────────────────────────────────────────
\begin{frame}{Skills: Responsabilidades e Outputs}
  \renewcommand{\arraystretch}{1.2}
  \begin{tabularx}{\textwidth}{@{}c p{2.8cm} X p{3.0cm}@{}}
    \toprule
    \textbf{\#} & \textbf{Nome} & \textbf{Responsabilidade} & \textbf{Output chave} \\
    \midrule
    01 & Escopo \& EAP & Normaliza serviços, cria itens da EAP com quantitativos e precedências & EAP validada, DAG de atividades \\
    02 & Prod.\ Climática & Consulta precipitação histórica; calcula $\lambda_j$ por semana ISO & \texttt{Praticabilidade} semanal \\
    03 & Catálogo CPU & Resolve CPUs, composições de equipe e PU de catálogo ($p_i^c$) & \texttt{EquipeDireta} por item \\
    04 & Cronograma & Resolve RCPSP: início/fim por atividade, $m_{ij}$, path crítico & \texttt{ScheduledActivity} + pool spec \\
    05 & Equipes \& Pool & Contrato de pool, custo instancial por lifecycle, $\tilde{C}_i$ & \texttt{EquipeDireta.custo\_total} \\
    06 & Resultado Financeiro & $p_i^o$, bridge N1$\to$N2$\to$N3, custo por categoria & \texttt{Planejamento}, \texttt{CustoTotal} \\
    07 & Indiretos \& BDI & TOP/LAB/ADM/QSMS/MOB + BDI por categoria; valor de venda & \texttt{Venda}, \texttt{MemorialBDI} \\
    08 & Artefatos & Empacota PDF memorial, XLSX, cubo analítico, benchmark ZIP & Bundle versionado \\
    \bottomrule
  \end{tabularx}
\end{frame}

% =============================================================================
\section{Regras de Negócio}
% =============================================================================

\begin{frame}{Regras de Negócio: Custo Operacional}
  \begin{columns}[T]
    \column{0.48\textwidth}
    \begin{block}{Custo horário de catálogo $\gamma_i$}
      \[
        \gamma_i = \sum_{r \in R_i} \gamma_r
        = \frac{\tilde{C}_i^{\text{sem}}}{h_{\text{sem}}}
      \]
      Decomposto em MO, EQ, combustível, material.
    \end{block}
    \vspace{4pt}
    \begin{block}{Preço unitário de catálogo $p_i^c$}
      \[
        p_i^c = \frac{\gamma_i}{\rho_i}
        = \frac{\sum_r C_r^{\text{sem}}}{P_i \cdot h_{\text{sem}}}
      \]
      $\rho_i$: produtividade horária da equipe-padrão.
    \end{block}
    \vspace{4pt}
    \begin{block}{Identidade de auditoria (NEP-043)}
      \[
        p_i^c \cdot Q_i = \tilde{C}_{i}^{\text{op}} \quad \text{se } \eta=1
      \]
      Verifica bridge skill-05 $\leftrightarrow$ skill-06.
    \end{block}

    \column{0.48\textwidth}
    \begin{block}{Preço unitário operacional $p_i^o$}
      \[
        p_i^o = p_i^c \cdot \underbrace{
          \frac{\sum_j m_{ij}}{\sum_j \delta_j\,\lambda_j\,m_{ij}}
        }_{\tau_i \;\geq\; 1}
      \]
      $\delta_j = d_j/d$: fração de dias na semana $j$.\\
      $\lambda_j$: praticabilidade climática.\\
      $\tau_i \geq 1$: custo operacional sempre $\geq$ catálogo.
    \end{block}
    \vspace{4pt}
    \begin{block}{Lifecycle engine (ADR-106)}
      \[
        \tilde{C}_i = \underbrace{c_{\text{mo}} \cdot n_{\text{prod}}}_{\text{ativo MO}}
        + \underbrace{c_{\text{eq}} \cdot n_{\text{prod}}}_{\text{ativo EQ}}
        + \underbrace{c_{\text{mo}} \cdot n_{\text{idle}}}_{\text{ocioso MO}}
        + \underbrace{c_{\text{eq}}^{\rho} \cdot n_{\text{idle}}}_{\text{ocioso EQ}}
        + c_{\text{mob/demob}}
      \]
    \end{block}
  \end{columns}
\end{frame}

\begin{frame}{Regras de Negócio: BDI e Valor de Venda}
  \begin{columns}[T]
    \column{0.48\textwidth}
    \begin{block}{Estrutura do BDI (ADR-023 / TCU)}
      Por categoria de recurso $r$:
      \[
        \beta_r = \frac{AC + DF + R + L + \text{tributos}}{1 - L - \text{tributos}}
      \]
      \medskip
      \begin{tabular}{ll}\scriptsize
        Material    & $\beta_{\text{mat}}$ \\
        Diesel      & $\beta_{\text{diesel}}$ \\
        Mão-de-obra & $\beta_{\text{mo}}$ \\
        Equipamentos & $\beta_{\text{eq}}$ \\
      \end{tabular}
    \end{block}
    \vspace{4pt}
    \begin{block}{Valor de venda direto}
      \[
        V_i^D = \sum_{j}\sum_{r \in A_i} \beta_r \cdot \gamma_r \cdot h \cdot d \cdot m_{ij}
      \]
    \end{block}

    \column{0.48\textwidth}
    \begin{block}{Valor de venda indireto}
      \[
        V_i^I = \sum_{j}\sum_{q \in B_i} \beta_r \cdot \gamma_q \cdot h \cdot d
      \]
      Indiretos (TOP/LAB/ADM/QSMS/MOB) não têm $m_{ij}$ --- permanência = duração do projeto.
    \end{block}
    \vspace{4pt}
    \begin{block}{Métricas gerenciais}
      \begin{align*}
        \frac{\tilde{C}^I}{\tilde{C}^D} &: \text{razão indireto/direto} \\[4pt]
        \frac{V - \tilde{C}}{\tilde{C}} &: \text{margem de operacionalização}
      \end{align*}
    \end{block}
    \vspace{4pt}
    \begin{alertblock}{\small Memorial BDI (ADR-057)}
      \scriptsize Trilha auditável: salário $\to$ custo $\to$ BDI $\to$ PV.
      Hash SHA-256 garante integridade.
    \end{alertblock}
  \end{columns}
\end{frame}

\begin{frame}{Regras de Negócio: Pool de Recursos (RCPSP)}
  \begin{columns}[T]
    \column{0.50\textwidth}
    \begin{block}{O que é o pool}
      Conjunto de recursos (equipamentos + equipes) compartilhados entre atividades:
      \[
        P \supseteq \biguplus_i R_i, \quad m_{ij} \leq \#\{R_i\}
      \]
      Contratos de capacidade (ADR-071): custo do pool é rateado proporcionalmente
      pela fração de atribuição $\phi_i$.
    \end{block}
    \vspace{4pt}
    \begin{block}{Scheduler RCPSP (skill-04)}
      Minimiza makespan sujeito a:
      \begin{itemize}\scriptsize
        \item Precedências FS/SS/FF com lag e $\theta$ de conclusão (ADR-081)
        \item Capacidade do pool (recursos simultâneos $\leq$ disponíveis)
        \item Janelas de praticabilidade $\lambda_j$
      \end{itemize}
      Solver: CP-SAT (OR-Tools) com fallback MINSLK.
    \end{block}

    \column{0.46\textwidth}
    \begin{block}{Cubo analítico de pool (ADR-086)}
      7 faces de análise:
      \begin{enumerate}\scriptsize
        \item \textbf{Ciclo de vida} de cada instância
        \item \textbf{Custo por estado} (ativo/ocioso/mob/demob)
        \item \textbf{Utilização temporal} (histograma semanal)
        \item \textbf{Atividade $\times$ instância}
        \item \textbf{Causa de custo} (produtivo/ocioso/mobilização)
        \item \textbf{Economia de pool} vs.\ dedicado
        \item \textbf{Cadeia de perpetuação} de custo
      \end{enumerate}
    \end{block}
    \vspace{4pt}
    \begin{exampleblock}{\small Campanha GO (what-if formal)}
      \scriptsize Executa $N$ cenários em paralelo variando regime,
      frentes e composição. Retorna shortlist Pareto $O \times EF$.
    \end{exampleblock}
  \end{columns}
\end{frame}

% =============================================================================
\section{Execução Versionada}
% =============================================================================

\begin{frame}{Execução Versionada: Escopos e Planos}
  \begin{columns}[T]
    \column{0.52\textwidth}
    \begin{block}{Hierarquia (ADR-031)}
      \begin{tikzpicture}[
          box/.style={draw=mlblue,fill=mlgray,rounded corners=2pt,
                      minimum width=3.0cm,font=\tiny,align=center,
                      inner sep=3pt},
          arr/.style={-{Stealth[length=3pt]},mlblue},
        ]
        \node[box,fill=mlblue!10] (proj) {Projeto};
        \node[box,below=0.2cm of proj] (esc) {Escopo (limites físicos)};
        \node[box,below=0.2cm of esc,fill=mlgreen!15] (plan_a) {Plano Principal \textbf{APROVADO}};
        \node[box,below=0.15cm of plan_a] (plan_b) {Plano Alternativo EM\_ANDAMENTO};
        \node[box,below=0.15cm of plan_b,fill=gray!15,text=gray] (plan_c) {Plano C \textit{ARQUIVADO}};
        \draw[arr] (proj) -- (esc);
        \draw[arr] (esc) -- (plan_a);
        \draw[arr,dashed] (plan_a) -- (plan_b);
        \draw[arr,dashed] (plan_b) -- (plan_c);
      \end{tikzpicture}
    \end{block}

    \column{0.44\textwidth}
    \begin{block}{Estados de plano}
      \[
        \texttt{RASCUNHO} \to \texttt{EM\_ANDAMENTO} \to \texttt{APROVADO}
      \]
      \[
        \texttt{EM\_ANDAMENTO} \to \texttt{ARQUIVADO}
      \]
      \vspace{2pt}
      {\scriptsize Apenas 1 plano APROVADO por escopo. Aprovação congela
      todos os artefatos e o snapshot do catálogo.}
    \end{block}
    \vspace{4pt}
    \begin{block}{O que cada plano preserva}
      \begin{itemize}\scriptsize
        \item Canônico JSON após cada skill
        \item Inputs de skill (overrides, pool\_config)
        \item Log de invocações e status
        \item Snapshot imutável do catálogo (ADR-041)
        \item Perfil BDI com hash SHA-256
        \item Artefatos PDF/XLSX/JSON
      \end{itemize}
    \end{block}
  \end{columns}
\end{frame}

\begin{frame}{Execução Versionada: Reprodutibilidade}
  \begin{columns}[T]
    \column{0.50\textwidth}
    \begin{block}{Snapshot de catálogo por execução}
      \begin{lstlisting}
projetos/{uuid}/snapshots/
  insumos_overlay.json   # precos usados
  cpus_bundle.json       # composicoes
  jurisdicao.json        # encargos/tributacao
      \end{lstlisting}
      \vspace{2pt}
      {\scriptsize Re-executar a pipeline com o mesmo snapshot
      produz o mesmo resultado --- mesmo que o catálogo tenha
      mudado depois.}
    \end{block}
    \vspace{4pt}
    \begin{block}{Rastreabilidade completa}
      \begin{lstlisting}
Projeto {uuid}
  Plano {plano_id}
    contexto_version: "v2"
    snapshot: v2_2026-06-01.json
    resultado.custo: R$ 43.8M
      \end{lstlisting}
    \end{block}

    \column{0.46\textwidth}
    \begin{block}{Jobs assíncronos (worker pool)}
      \begin{tikzpicture}[
          st/.style={draw=mlblue,fill=mlgray,rounded corners=2pt,
                     minimum width=1.6cm,font=\tiny,align=center,
                     inner sep=3pt},
          arr/.style={-{Stealth[length=3pt]},mlblue,thick},
        ]
        \node[st] (q) {AGUARDANDO\\FILA};
        \node[st,right=0.3cm of q] (r) {EM\\EXECUCAO};
        \node[st,fill=mlgreen!20,above right=0.2cm and 0.3cm of r] (ok) {CONCLUIDO};
        \node[st,fill=mlred!20,below right=0.2cm and 0.3cm of r] (err) {ERRO};
        \node[st,fill=gray!20,right=0.3cm of ok] (can) {CANCELADO};
        \draw[arr] (q) -- (r);
        \draw[arr] (r) -- (ok);
        \draw[arr] (r) -- (err);
        \draw[arr] (r) -- (can);
      \end{tikzpicture}
      \vspace{2pt}
      {\scriptsize \texttt{disparar\_pipeline} retorna \texttt{job\_id}
      imediatamente. Polling via \texttt{consultar\_job\_pipeline}.}
    \end{block}
    \vspace{4pt}
    \begin{alertblock}{\small Comparação entre planos}
      \scriptsize \texttt{comparar\_planos(plano\_a, plano\_b)} retorna
      lado a lado: custo direto/indireto, margem, prazo, composição BDI.
    \end{alertblock}
    \vspace{4pt}
    \begin{block}{BDI versionado dentro do projeto}
      \scriptsize \texttt{salvar\_perfil\_bdi} gera \texttt{perfil\_id},
      versão incremental e diff campo a campo.
      Hash de conteúdo previne adulteração silenciosa.
    \end{block}
  \end{columns}
\end{frame}

% =============================================================================
\section{Gerenciamento de Catálogo}
% =============================================================================

\begin{frame}{Catálogo: Estrutura e Versionamento (ADR-062)}
  \begin{columns}[T]
    \column{0.50\textwidth}
    \begin{block}{O que compõe um contexto}
      \begin{itemize}\scriptsize
        \item \textbf{Insumos} --- preços de MO, EQ, material, diesel
        \item \textbf{CPUs} --- composições padrão de equipe + PU
        \item \textbf{Produções} --- $\rho_i$ por condição/variante
        \item \textbf{Jurisdições} --- encargos e regime tributário
      \end{itemize}
    \end{block}
    \vspace{4pt}
    \begin{block}{Estrutura de arquivos}
      \begin{lstlisting}
config/contextos/mendeslobo/
  contexto.json      # metadados
  current            # "v3"
  versions/
    v1.json
    v2.json          # SINAPI jun/26
    v3.json          # diesel set/26
      \end{lstlisting}
    \end{block}
    \vspace{4pt}
    \begin{alertblock}{\small Isolamento de projetos}
      \scriptsize Projetos existentes \textbf{não são afetados} por nova versão
      de catálogo. Cada projeto captura \texttt{contexto\_version} na criação.
    \end{alertblock}

    \column{0.46\textwidth}
    \begin{block}{Fluxo de atualização}
      \begin{tikzpicture}[
          box/.style={draw=mlblue,fill=mlgray,rounded corners=2pt,
                      minimum width=3.4cm,font=\tiny,align=center,
                      inner sep=3pt},
          arr/.style={-{Stealth[length=3pt]},mlblue},
        ]
        \node[box] (rfi) {Pesquisa de mercado / SINAPI};
        \node[box,below=0.2cm of rfi] (cmp) {\texttt{comparar\_rfi\_catalogo}\\divergências vs.\ master};
        \node[box,below=0.2cm of cmp] (upd) {\texttt{atualizar\_preco\_insumo}\\rastreabilidade: quem, quando, motivo};
        \node[box,below=0.2cm of upd] (ras) {\texttt{criar\_rascunho\_catalogo}};
        \node[box,below=0.2cm of ras,fill=mlgreen!15] (mrg) {\texttt{aprovar\_merge\_catalogo}\\(admin / cat-manager)};
        \draw[arr] (rfi) -- (cmp);
        \draw[arr] (cmp) -- (upd);
        \draw[arr] (upd) -- (ras);
        \draw[arr] (ras) -- (mrg);
      \end{tikzpicture}
    \end{block}
    \vspace{4pt}
    \begin{exampleblock}{\small Regra dos quatro olhos}
      \scriptsize \texttt{engineer} pode criar e editar rascunho, mas
      \textbf{não pode} fazer merge. Apenas \texttt{admin} ou \texttt{cat-manager}
      aprovam o merge.
    \end{exampleblock}
  \end{columns}
\end{frame}

\begin{frame}{Catálogo: Diagnóstico e Previsão de Impacto}
  \begin{columns}[T]
    \column{0.50\textwidth}
    \begin{block}{Ferramentas de diagnóstico}
      \begin{tabularx}{\linewidth}{@{}l X@{}}
        \toprule
        \textbf{Tool} & \textbf{O que faz} \\
        \midrule
        \tool{auditar\_catalogo} & \% CPUs calculáveis, inconsistências \\
        \tool{listar\_divergencias} & Diff entre rascunho e master \\
        \tool{consultar\_impacto} & $\Delta$PU em todos os itens afetados \\
        \tool{historico\_preco} & Série histórica de insumo \\
        \tool{calcular\_tributos} & Encargos CLT/PJ sem persistir \\
        \bottomrule
      \end{tabularx}
    \end{block}
    \vspace{4pt}
    \begin{block}{Simulação sem persistir}
      \begin{lstlisting}
calcular_pu_producao(
  contexto="mendeslobo",
  cpu_id="PROD-TERRA-ESCO-001",
  preco_diesel=6.80
)
# Retorna PU recalculada -- sem alterar catálogo
      \end{lstlisting}
    \end{block}

    \column{0.46\textwidth}
    \begin{block}{Produtividade como artefato (ADR-007)}
      \begin{itemize}\scriptsize
        \item \texttt{Producao.valor\_horario} --- $\rho_i$ versionado
        \item Origem: calibrado contra obra real, tabela de referência ou regressão
        \item Qualidade: intervalo de confiança, $n$ amostras, desvio padrão
        \item \texttt{fator\_producao(nova\_composicao)}: escala $\rho$ para equipes alternativas
      \end{itemize}
    \end{block}
    \vspace{4pt}
    \begin{block}{Encargos e regime tributário (ADR-097)}
      \begin{tabular}{ll}\scriptsize
        Horista CLT  & 177{,}70\% \quad (vigente abr/2026) \\
        Mensalista   & 106{,}00\% \\
        PJ           & 6{,}00\% \\
      \end{tabular}\\[4pt]
      {\scriptsize Grupos A--F separados para reconciliação fe\_nec com $PU_{\text{sbdi}}$ (NEP-037).}
    \end{block}
    \vspace{4pt}
    \begin{alertblock}{\small Benchmark (NEP-039)}
      \scriptsize Projetos de referência são \textbf{imutáveis} após aprovação.
      Edição cria versão $v+1$ automaticamente.
    \end{alertblock}
  \end{columns}
\end{frame}

% =============================================================================
\section{Roles e Tools}
% =============================================================================

\begin{frame}{Roles: Cinco Perfis de Acesso}
  \begin{columns}[T]
    \column{0.56\textwidth}
    \begin{block}{Perfis e responsabilidades}
      \renewcommand{\arraystretch}{1.15}
      \begin{tabularx}{\linewidth}{@{}l X r@{}}
        \toprule
        \textbf{Role} & \textbf{Responsabilidade principal} & \textbf{Tools} \\
        \midrule
        \role{operator}     & Precificação RFI baseada em catálogo & 36 \\
        \role{engineer}     & Planejamento completo RFI + RFP      & 40 \\
        \role{admin}        & Aprovações e gestão de catálogo       & 42 \\
        \role{auditor}      & Leitura de projetos e artefatos       & 16 \\
        \role{cat-manager}  & Gestão exclusiva do catálogo          & 18 \\
        \bottomrule
      \end{tabularx}
    \end{block}
    \vspace{4pt}
    \begin{block}{Permissões granulares (ADR-099)}
      \begin{tabular}{ll}\scriptsize
        \texttt{project:read}    & Consulta de projetos/artefatos \\
        \texttt{project:create}  & Criação e edição de projetos \\
        \texttt{project:run}     & Execução da pipeline \\
        \texttt{project:approve} & Aprovação de versões \\
        \texttt{catalog:read}    & Leitura do catálogo \\
        \texttt{catalog:write}   & Edição de catálogo \\
        \texttt{catalog:merge}   & Aprovação de merge \\
        \texttt{artifacts:*}     & Leitura e download \\
        \texttt{financials:read} & Resultados de campanha \\
      \end{tabular}
    \end{block}

    \column{0.40\textwidth}
    \begin{alertblock}{\small Intent RFI vs.\ RFP}
      \scriptsize \role{operator} só pode disparar pipeline com
      \texttt{intent=RFI} (estimativa de preço).\\[4pt]
      \role{engineer}/\role{admin} podem disparar com
      \texttt{intent=RFP} (planejamento completo com pool e cronograma detalhado).\\[4pt]
      Validado em \texttt{permission\_guard.py} via \texttt{INTENT\_PERMISSIONS}.
    \end{alertblock}
    \vspace{4pt}
    \begin{exampleblock}{\small Cat-manager sem projetos}
      \scriptsize \role{cat-manager} não tem acesso a projetos ---
      vê apenas catálogo e resultados de campanha.
      Segregação de domínio intencional.
    \end{exampleblock}
  \end{columns}
\end{frame}

\begin{frame}{Matriz Role × Tool (resumo)}
  \setlength{\tabcolsep}{4pt}
  \renewcommand{\arraystretch}{1.1}
  \begin{tabularx}{\textwidth}{@{}X c c c c c@{}}
    \toprule
    \textbf{Grupo / Tool} & \role{oper} & \role{eng} & \role{adm} & \role{aud} & \role{cat} \\
    \midrule
    \multicolumn{6}{@{}l}{\scriptsize\textit{Sistema}} \\
    \tool{ping}                      & \yes & \yes & \yes & \yes & \yes \\[2pt]
    \multicolumn{6}{@{}l}{\scriptsize\textit{Projetos --- leitura}} \\
    \tool{listar\_projetos}, \tool{consultar\_status}, \tool{comparar\_planos}\ldots & \yes & \yes & \yes & \yes & \no \\[2pt]
    \multicolumn{6}{@{}l}{\scriptsize\textit{Projetos --- criação}} \\
    \tool{criar\_projeto}, \tool{iniciar\_orcamento}, \tool{nova\_versao}\ldots & \yes & \yes & \yes & \no & \no \\[2pt]
    \multicolumn{6}{@{}l}{\scriptsize\textit{Projetos --- execução}} \\
    \tool{disparar\_pipeline}        & \half~RFI & \yes & \yes & \no & \no \\
    \tool{explorar\_configuracoes\_go} & \half~RFI & \yes & \yes & \no & \no \\
    \tool{simular\_cenario\_replan}  & \yes & \yes & \yes & \no & \no \\[2pt]
    \multicolumn{6}{@{}l}{\scriptsize\textit{Projetos --- aprovação}} \\
    \tool{aprovar\_versao}           & \no & \yes & \yes & \no & \no \\[2pt]
    \multicolumn{6}{@{}l}{\scriptsize\textit{Catálogo --- leitura}} \\
    \tool{pesquisar\_catalogo}, \tool{historico\_preco}\ldots & \yes & \yes & \yes & \no & \yes \\[2pt]
    \multicolumn{6}{@{}l}{\scriptsize\textit{Catálogo --- escrita}} \\
    \tool{atualizar\_preco\_insumo}, \tool{criar\_rascunho}\ldots & \no & \yes & \yes & \no & \yes \\[2pt]
    \multicolumn{6}{@{}l}{\scriptsize\textit{Catálogo --- merge}} \\
    \tool{aprovar\_merge\_catalogo}  & \no & \no & \yes & \no & \yes \\[2pt]
    \multicolumn{6}{@{}l}{\scriptsize\textit{Financeiro}} \\
    \tool{obter\_resultado\_campanha}, \tool{comparar\_rfi\_catalogo} & \yes & \yes & \yes & \yes & \yes \\
    \bottomrule
    \multicolumn{6}{@{}l}{\scriptsize \yes~acesso pleno \quad \half~acesso parcial (RFI) \quad \no~sem permissão} \\
  \end{tabularx}
\end{frame}

% =============================================================================
\section{Workflows por Role}
% =============================================================================

\begin{frame}{Workflow: Operator}
  \begin{columns}[T]
    \column{0.52\textwidth}
    \begin{block}{Caso de uso principal: estimativa de preço (RFI)}
      \begin{tikzpicture}[
          box/.style={draw=mlblue,fill=mlgray,rounded corners=2pt,
                      minimum width=4.0cm,font=\scriptsize,
                      align=left,inner sep=4pt},
          arr/.style={-{Stealth[length=3.5pt]},mlblue,thick},
        ]
        \node[box] (p1) {\texttt{ping()} --- confirmar sistema OK};
        \node[box,below=0.18cm of p1] (p2) {\texttt{listar\_contextos()} + \texttt{pesquisar\_catalogo()}};
        \node[box,below=0.18cm of p2] (p3) {\texttt{criar\_projeto\_formulario()} ou \texttt{iniciar\_orcamento()}};
        \node[box,below=0.18cm of p3] (p4) {\texttt{disparar\_pipeline(intent=\textbf{RFI})}};
        \node[box,below=0.18cm of p4] (p5) {\texttt{consultar\_job\_pipeline()} --- aguardar};
        \node[box,below=0.18cm of p5] (p6) {\texttt{consultar\_status()} + \texttt{auditar\_plano()}};
        \node[box,fill=mlorange!20,below=0.18cm of p6] (p7) {\faArrowRight\ \textbf{HANDOFF} para engineer};
        \foreach \a/\b in {p1/p2,p2/p3,p3/p4,p4/p5,p5/p6,p6/p7}
          \draw[arr] (\a) -- (\b);
      \end{tikzpicture}
    \end{block}

    \column{0.44\textwidth}
    \begin{block}{O que operator pode e não pode}
      \textbf{Pode:}
      \begin{itemize}\scriptsize
        \item Criar projetos e iniciar orçamentos
        \item Executar pipeline completa com \texttt{intent=RFI}
        \item Consultar catálogo, histórico de preços
        \item Simular cenários lineares (\texttt{simular\_cenario\_replan})
        \item Acessar artefatos e benchmarks
        \item Empacotar benchmark ZIP
      \end{itemize}
      \vspace{4pt}
      \textbf{Não pode:}
      \begin{itemize}\scriptsize
        \item Disparar pipeline com \texttt{intent=RFP}
        \item Aprovar versões de planos
        \item Editar catálogo
      \end{itemize}
    \end{block}
    \vspace{4pt}
    \begin{alertblock}{\small Handoff H1}
      \scriptsize Após \texttt{disparar\_pipeline(RFI)}, a versão entra em
      \texttt{EM\_VALIDACAO}. Engineer ou admin executam \texttt{aprovar\_versao}.
    \end{alertblock}
  \end{columns}
\end{frame}

\begin{frame}{Workflow: Engineer}
  \begin{columns}[T]
    \column{0.52\textwidth}
    \begin{block}{Caso de uso: planejamento completo (RFP)}
      \begin{tikzpicture}[
          box/.style={draw=mlblue,fill=mlgray,rounded corners=2pt,
                      minimum width=4.2cm,font=\scriptsize,
                      align=left,inner sep=4pt},
          arr/.style={-{Stealth[length=3.5pt]},mlblue,thick},
        ]
        \node[box] (e1) {Skills 01--03 via \texttt{disparar\_pipeline(RFI)}};
        \node[box,below=0.18cm of e1] (e2) {\texttt{explorar\_configuracoes\_go(RFP)} --- sweep};
        \node[box,below=0.18cm of e2] (e3) {\texttt{disparar\_pipeline(intent=\textbf{RFP})} --- skills 04--08};
        \node[box,below=0.18cm of e3] (e4) {\texttt{explorar\_pool\_recursos()} --- 7 faces analíticas};
        \node[box,below=0.18cm of e4] (e5) {\texttt{comparar\_planos()} --- \texttt{nova\_versao()} se necessário};
        \node[box,fill=mlgreen!20,below=0.18cm of e5] (e6) {\texttt{aprovar\_versao()} --- plano congelado};
        \node[box,fill=mlorange!20,below=0.18cm of e6] (e7) {\faArrowRight\ Rascunho catálogo \textbf{HANDOFF} admin};
        \foreach \a/\b in {e1/e2,e2/e3,e3/e4,e4/e5,e5/e6,e6/e7}
          \draw[arr] (\a) -- (\b);
      \end{tikzpicture}
    \end{block}

    \column{0.44\textwidth}
    \begin{block}{Capacidades exclusivas}
      \begin{itemize}\scriptsize
        \item \texttt{aprovar\_versao} (operator não pode)
        \item \texttt{disparar\_pipeline(intent=RFP)}
        \item Editar catálogo (\texttt{atualizar\_preco\_insumo},\\
              \texttt{adicionar\_variante\_producao})
        \item Criar rascunho de catálogo
        \item \texttt{salvar\_perfil\_bdi} com versionamento
      \end{itemize}
    \end{block}
    \vspace{4pt}
    \begin{block}{Handoffs gerados pelo engineer}
      \textbf{H2} Exploração RFI $\to$ análise RFP completa\\[2pt]
      \textbf{H3} Rascunho catálogo $\to$ admin/cat-manager para merge
      \begin{lstlisting}
engineer:
  criar_rascunho_catalogo()
  atualizar_preco_insumo()
  listar_divergencias_rascunho()
    -- HANDOFF --
  admin / cat-manager:
    aprovar_merge_catalogo()
      \end{lstlisting}
    \end{block}
  \end{columns}
\end{frame}

\begin{frame}{Workflows: Admin, Auditor e Cat-Manager}
  \begin{columns}[T]
    \column{0.32\textwidth}
    \begin{block}{\role{admin} (42 tools)}
      \begin{itemize}\scriptsize
        \item Tudo de engineer +
        \item \texttt{aprovar\_merge\_catalogo}
        \item \texttt{rejeitar\_rascunho\_catalogo}
        \item Gestão de usuários (CLI)
        \item Aprovação de versões e merges de catálogo sem restrição
      \end{itemize}
      \vspace{4pt}
      \textbf{Handoff H4:}\\
      {\scriptsize CLI cria usuário $\to$ email com \texttt{claude\_desktop\_config.json} $\to$ usuário inicia sessão MCP}
    \end{block}

    \column{0.32\textwidth}
    \begin{block}{\role{auditor} (16 tools)}
      \begin{itemize}\scriptsize
        \item Somente leitura de projetos e artefatos
        \item Acessa resultado de campanha (\texttt{obter\_resultado\_campanha})
        \item Compara RFI vs.\ catálogo (\texttt{comparar\_rfi\_catalogo})
        \item \textbf{Não} cria, edita, executa ou aprova nada
        \item Caso de uso: auditoria de conformidade, análise pós-obra
      \end{itemize}
      \vspace{4pt}
      {\scriptsize \texttt{auditar\_plano} e \texttt{download\_benchmark\_zip} são as
      principais ferramentas de trabalho.}
    \end{block}

    \column{0.32\textwidth}
    \begin{block}{\role{cat-manager} (18 tools)}
      \begin{itemize}\scriptsize
        \item Gestão completa do catálogo
        \item Cria rascunho, edita insumos, aprova merge
        \item Acessa \texttt{obter\_resultado\_campanha} e \texttt{comparar\_rfi\_catalogo}
        \item \textbf{Não} vê projetos, planos ou artefatos de obra
      \end{itemize}
      \vspace{4pt}
      \textbf{Workflow típico:}\\
      {\scriptsize \texttt{criar\_rascunho} $\to$ \texttt{atualizar\_preco\_insumo}  $\to$ \texttt{listar\_divergencias} $\to$ \texttt{aprovar\_merge\_catalogo}}
    \end{block}
  \end{columns}
\end{frame}

\begin{frame}{Handoffs entre Roles}
  \begin{tikzpicture}[
      role/.style={draw=mlblue,fill=mlblue!10,rounded corners=3pt,
                   minimum width=2.0cm,minimum height=0.7cm,
                   font=\small\bfseries,align=center},
      act/.style={draw=mlblue!50,fill=mlgray,rounded corners=2pt,
                  minimum width=3.4cm,font=\scriptsize,align=center,
                  inner sep=3pt},
      arr/.style={-{Stealth[length=4pt]},mlblue,thick},
      hoff/.style={-{Stealth[length=4pt]},mlorange,thick,dashed},
    ]
    % H1
    \node[role] (op1) {operator};
    \node[act,right=0.4cm of op1]  (a1) {\texttt{disparar\_pipeline(RFI)}\\versão EM\_VALIDACAO};
    \node[role,right=0.4cm of a1]  (en1) {engineer};
    \node[act,right=0.4cm of en1]  (b1) {\texttt{aprovar\_versao()}\\versão APROVADO};
    \draw[arr] (op1) -- (a1);
    \draw[hoff] (a1) -- node[above,font=\tiny\itshape,mlorange]{H1} (en1);
    \draw[arr] (en1) -- (b1);

    % H2
    \node[role,below=0.7cm of op1] (op2) {operator};
    \node[act,right=0.4cm of op2]  (a2) {\texttt{explorar\_go(RFI)}\\campanha inicial};
    \node[role,right=0.4cm of a2]  (en2) {engineer};
    \node[act,right=0.4cm of en2]  (b2) {\texttt{explorar\_go(RFP)}\\planejamento completo};
    \draw[arr] (op2) -- (a2);
    \draw[hoff] (a2) -- node[above,font=\tiny\itshape,mlorange]{H2} (en2);
    \draw[arr] (en2) -- (b2);

    % H3
    \node[role,below=0.7cm of op2] (en3) {engineer};
    \node[act,right=0.4cm of en3]  (a3) {\texttt{criar\_rascunho}\\+ \texttt{atualizar\_preco}};
    \node[role,right=0.4cm of a3]  (ad3) {admin /\\cat-mgr};
    \node[act,right=0.4cm of ad3]  (b3) {\texttt{aprovar\_merge}\\catálogo publicado};
    \draw[arr] (en3) -- (a3);
    \draw[hoff] (a3) -- node[above,font=\tiny\itshape,mlorange]{H3} (ad3);
    \draw[arr] (ad3) -- (b3);

    % H4
    \node[role,below=0.7cm of en3] (ad4) {admin};
    \node[act,right=0.4cm of ad4]  (a4) {CLI: criar usuário\\email config};
    \node[role,right=0.4cm of a4]  (us4) {novo\\usuário};
    \node[act,right=0.4cm of us4]  (b4) {configura Claude Desktop\\inicia sessão MCP};
    \draw[arr] (ad4) -- (a4);
    \draw[hoff] (a4) -- node[above,font=\tiny\itshape,mlorange]{H4} (us4);
    \draw[arr] (us4) -- (b4);
  \end{tikzpicture}
  \vspace{6pt}
  {\scriptsize
    \textbf{H1} Execução $\to$ aprovação (operator $\to$ engineer/admin) \quad
    \textbf{H2} RFI $\to$ RFP completo (operator $\to$ engineer) \\
    \textbf{H3} Rascunho catálogo $\to$ merge (engineer $\to$ admin/cat-manager) \quad
    \textbf{H4} Onboarding (admin $\to$ novo usuário)
  }
\end{frame}

% =============================================================================
\section{Infraestrutura}
% =============================================================================

\begin{frame}{Infraestrutura: Azure ACA, Blue-Green e CI/CD}
  \begin{columns}[T]
    \column{0.50\textwidth}
    \begin{block}{Azure Container Apps (ACA)}
      \begin{itemize}\scriptsize
        \item Migrado de ACI para ACA (blue-green nativo via revisions)
        \item Blue-green sem downtime: nova revisão recebe 0\% tráfego,
              smoke pass $\to$ 100\% tráfego, revisão anterior desmobilizada
        \item 14 checks de smoke test (saúde, catálogo, projetos, worker pool)
      \end{itemize}
    \end{block}
    \vspace{4pt}
    \begin{block}{Pipeline CI/CD (GitHub Actions)}
      \begin{tabular}{ll}\scriptsize
        \texttt{ci.yml}     & lint + typecheck + testes unitários \\
        \texttt{deploy.yml} & blue-green deploy via \texttt{deploy.sh} \\
        \texttt{infra.yml}  & Bicep --- infraestrutura como código \\
      \end{tabular}\\[4pt]
      {\scriptsize \texttt{deploy.sh} é o entrypoint único: image / deploy /
      smoke / promote / rollback / status / bootstrap / keys}
    \end{block}
    \vspace{4pt}
    \begin{alertblock}{\small Rollback automático}
      \scriptsize Se smoke falha, rollback automático restitui revisão anterior.
      Sem intervenção manual.
    \end{alertblock}

    \column{0.46\textwidth}
    \begin{block}{Segurança e autenticação (ADR-099)}
      \begin{itemize}\scriptsize
        \item JWT obrigatório em todas as tools MCP
        \item \texttt{McpAuthMiddleware} valida token e role
        \item \texttt{PermissionGuard} verifica permissão por tool
        \item \texttt{INTENT\_PERMISSIONS} para restrição de intent
        \item Rate limits por operação crítica
      \end{itemize}
    \end{block}
    \vspace{4pt}
    \begin{block}{Observabilidade}
      \begin{itemize}\scriptsize
        \item OpenTelemetry (LoggingHandler + traces)
        \item Logfire desabilitável via \texttt{PYDANTIC\_DISABLE\_PLUGINS}
        \item Notificações de resultado via email (ADR-123)
        \item Worker pool com 4 workers e fila de jobs
      \end{itemize}
    \end{block}
    \vspace{4pt}
    \begin{exampleblock}{\small SharePoint (em implantação)}
      \scriptsize Integração SharePoint para distribuição de artefatos
      e relatórios aos stakeholders.
    \end{exampleblock}
  \end{columns}
\end{frame}

% =============================================================================
\section{Encerramento}
% =============================================================================

\begin{frame}{Resumo Executivo}
  \begin{columns}[T]
    \column{0.48\textwidth}
    \begin{block}{O que entregamos}
      \begin{enumerate}\small
        \item \textbf{Pipeline auditável} de 8 skills: escopo $\to$ cronograma $\to$ pool $\to$ resultado financeiro
        \item \textbf{Versionamento completo}: planos, catálogo, BDI, snapshots imutáveis
        \item \textbf{Regras de negócio formalizadas}: $p_i^c$, $p_i^o = p_i^c \cdot \tau_i$, lifecycle engine, BDI TCU
        \item \textbf{5 roles} com 43 tools e regra dos quatro olhos para catálogo
        \item \textbf{Infraestrutura blue-green} com zero downtime e rollback automático
      \end{enumerate}
    \end{block}

    \column{0.48\textwidth}
    \begin{block}{Garantias do sistema}
      \begin{itemize}\small
        \item \textbf{Reprodutibilidade}: snapshot por execução garante mesmo resultado com mesmo input
        \item \textbf{Auditabilidade}: trilha completa catálogo $\to$ custo $\to$ BDI $\to$ venda com hash
        \item \textbf{Isolamento}: mudança de catálogo não quebra projetos existentes
        \item \textbf{Otimização}: RCPSP com CP-SAT minimiza makespan respeitando pool
        \item \textbf{What-if}: campanha GO com $N$ cenários Pareto $O \times EF$
      \end{itemize}
    \end{block}
  \end{columns}
  \vspace{8pt}
  \begin{center}
    \Large\textbf{Concilium} \\[4pt]
    {\normalsize Orçamentação de infraestrutura: auditável, versionada, otimizada.}
  \end{center}
\end{frame}

\end{document}
